\chapter{ماذا تفعل العصبونات في النوم}
\chaptermark{النوم}
\label{sec:sleep}
\begin{chaptergoals}
\item لماذا النوم مجموعة أوضاع يبدّلها \nt{ACET} و\nt{NORA} و\nt{SERO}، لا إطفاء؛
\item كيف يكوّن ضغط النوم وعتبتان مُرحِّلًا بتخلّف يضبطه الإيقاع اليومي؛
\item كيف يحوّل التعب القشرة إلى مذبذب استرخائي، وكيف يوقف \nt{ACET} الأرجحة؛
\item كيف يُعاد تشغيل اليوم مسرَّعًا، ولماذا قد تضعف الأوزان كلها بالنسبة نفسها في النوم.
\end{chaptergoals}

\section{النوم ليس إطفاءً، بل وضع آخر}

المهندس الذي يرى حاسوبًا مطفأً سيقول: إنه لا يفعل شيئًا. أما الدماغ النائم فليس كذلك. العصبونات لا تصمت في النوم: في النوم العميق تُصدر القشرة نبضات، وفي النوم السريع تُصدر منها قرابة ما تُصدره نهارًا. الذي يتغير ليس \emph{هل} تعمل العصبونات، بل \emph{كيف} تعمل. النوم مجموعة من أوضاع الآلة، وتبدّلها قنوات البث نفسها التي تعمل نهارًا.

لكل وضع يمكن كتابة ”حالة السجلات“ (الجدول~\ref{tab:sleepmodes}). في النهار تعمل \nt{ACET} و\nt{NORA} و\nt{SERO}، والمستشعرات موصولة، والعضلات تطيع. في النوم البطيء العميق تهدأ القنوات الثلاث كلها، ويضعف المدخل من أعضاء الحس~\cite{Hasselmo1999}: يهبط إطلاق الأسيتيل كولين في القشرة، وتُصدر عصبونات \nt{NORA} و\nt{SERO} نبضات أقل مما في النهار~\cite{Marrosu1995,AstonJonesBloom1981,McGintyHarper1976}. وفي النوم السريع صورة غريبة: يرتفع \nt{ACET} من جديد إلى مستوى النهار~\cite{Marrosu1995}، بينما يصمت \nt{NORA} و\nt{SERO} صمتًا شبه تام~\cite{AstonJonesBloom1981,McGintyHarper1976,JacobsAzmitia1992}. وعضلات الجسم مطفأة في النوم السريع: عصبونات \nt{ACET} التي تتحكم فيها تُثبَّط بقوة عبر \nt{GABA} و\nt{GLYC}~\cite{BrooksPeever2012}، بالمنع نفسه الذي رأيناه في الفصل~\ref{sec:gaba}، لكنه مفروض على المخرج كله دفعة واحدة.

\begin{table}[t]
\centering
\footnotesize
\setlength{\tabcolsep}{4pt}
\renewcommand{\arraystretch}{1.15}
\begin{tabular}{>{\raggedright}p{3.1cm}>{\raggedright}p{2.9cm}>{\raggedright}p{3.2cm}>{\raggedright\arraybackslash}p{3.4cm}}
\toprule
& اليقظة & النوم البطيء & النوم السريع \\
\midrule
\nt{ACET} (الكتابة، الفصل~\ref{sec:acet}) & مرتفع & منخفض & مرتفع \\
\nt{NORA} (الكسب، الفصل~\ref{sec:nora}) & متوسط، مع ومضات & منخفض & شبه صامت \\
\nt{SERO} (الفصل~\ref{sec:serocomputer}) & متوسط & منخفض & شبه صامت \\
المدخل من المستشعرات & موصول & مُضعَف & مُضعَف \\
المخرج إلى العضلات & موصول & مُضعَف & مطفأ بالتثبيط \\
نشاط القشرة & منتظم & أرجحة بطيئة: عمل، صمت & منتظم، كما في النهار \\
\bottomrule
\end{tabular}
\caption{حالة سجلات الآلة في الأوضاع الثلاثة. كل الأسطر مُقاسة؛ ومستويات \nt{ACET} و\nt{NORA} و\nt{SERO} بتسجيلات مباشرة بحسب مراحل النوم~\cite{Marrosu1995,AstonJonesBloom1981,McGintyHarper1976}.}
\label{tab:sleepmodes}
\end{table}

\section{متى ننام: مُرحِّل بتخلّف}

ما الذي يقرر متى تنتقل الآلة إلى النوم؟ يعمل هنا جيدًا مخطط بسيط من جزأين~\cite{Borbely1982}. الجزء الأول \emph{ضغط النوم} $S$، يتراكم ما دمنا مستيقظين ويُفرَّغ ما دمنا نائمين. إنه مكثف يُشحن نهارًا ويُفرَّغ ليلًا:
\begin{empheq}[box=\keybox]{equation}
\frac{\dd S}{\dd t} \;=\; \frac{1-S}{\tau_r}\ \ \text{(اليقظة)},\qquad
\frac{\dd S}{\dd t} \;=\; -\frac{S}{\tau_d}\ \ \text{(النوم)} .
\label{eq:sleeppressure}
\end{empheq}
الجزء الثاني عتبتان: تنام الآلة حين يبلغ $S$ العتبة العليا $H(t)$، وتستيقظ حين ينخفض $S$ إلى الدنيا $L(t)$. وكلتا العتبتين تتأرجح بسلاسة مع الإيقاع اليومي من الفصل~\ref{sec:chemoutput}. فنحصل على مُرحِّل بتخلّف، أي قادح شميت من القسم~\ref{sec:bistable}، ويعطي مع المكثف مذبذبًا استرخائيًا يضبط الإيقاعُ اليومي طورَه~\eqref{eq:pll}.

\begin{figure}[t]
\centering
\begin{tikzpicture}[>=Stealth,x=0.16cm,y=4.2cm]
\foreach \a/\b in {0/4.3,21.2/29.3,45.4/53.4,69.5/72}{\fill[blue!8] (\a,0) rectangle (\b,1);}
\draw[->,gray!70!black] (0,0) -- (75,0) node[right,font=\small,black] {$t$، ساعة};
\draw[->,gray!70!black] (0,0) -- (0,1.08) node[left,font=\small,black] {$S$};
\foreach \t in {0,24,48,72}{\draw[gray!60!black] (\t,-0.02) -- (\t,0.02); \node[below,font=\scriptsize] at (\t,0) {\t};}
\draw[thick,densely dashed,red!70!black] plot file {figures/sleep_H.dat};
\draw[thick,densely dashed,green!45!black] plot file {figures/sleep_L.dat};
\draw[very thick,blue!60!black] plot file {figures/sleep_S.dat};
\node[font=\scriptsize,red!70!black,anchor=west] at (73,0.72) {$H$: النوم};
\node[font=\scriptsize,green!45!black,anchor=west] at (73,0.24) {$L$: الاستيقاظ};
\node[font=\scriptsize,blue!60!black] at (25.25,0.93) {نوم};
\node[font=\scriptsize,blue!60!black] at (49.4,0.93) {نوم};
\end{tikzpicture}
\caption{ضغط النوم $S$~\eqref{eq:sleeppressure} عند $\tau_r=18.2$ ساعة و$\tau_d=4.2$ ساعة. في النهار يُشحن $S$ حتى العتبة العليا $H$، وفي الليل (المظلَّل) يُفرَّغ حتى الدنيا $L$؛ والعتبتان تتأرجحان مع الإيقاع اليومي. تصل الآلة بنفسها إلى نظام من نحو $16$ ساعة يقظة و$8$ ساعات نوم.}
\label{fig:twoprocess}
\end{figure}

في نموذجنا، بالقيم المعتادة $\tau_r=18.2$ ساعة و$\tau_d=4.2$ ساعة، تصل الآلة بنفسها إلى $16.1$ ساعة يقظة و$8.0$ ساعات نوم (الشكل~\ref{fig:twoprocess}). ويفسر النموذج أيضًا ما يبدو غريبًا: بعد أربعين ساعة بلا نوم يبلغ الضغط $0.90$، لكن نوم التعافي في النموذج لا يدوم إلا $8.8$ ساعة، لا يومًا كاملًا زيادة على المعتاد. فالتفريغ أُسّي، والشحنة العالية تزول بسرعة؛ و”الدَّين“ يُسدَّد بعمق النوم لا بطوله.

ما الذي يؤدي دور $S$ فيزيائيًا؟ المرشح القوي هو الأدينوزين، ”عدّاد الحمل“ من الملحق~\ref{sec:othermediators}: يرتفع تركيزه خلال اليقظة وينخفض في النوم~\cite{PorkkaHeiskanen1997,Dunwiddie2001}. أما أن ضغط النوم هو الأدينوزين بالذات ولا شيء غيره ففرضية.

\section{النوم البطيء: الشبكة كلها مذبذب استرخائي}

في النوم العميق تعمل عصبونات القشرة بالتناوب مجموعةً كاملة: أقل من مرة في الثانية تشتغل معًا لبضع مئات من أجزاء الألف من الثانية (\emph{حالة العمل}) وتصمت معًا (\emph{حالة الصمت}): تردد هذه الأرجحة أقل من $1\,\text{Hz}$، وتحت التخدير غالبًا $0.3\text{--}0.4\,\text{Hz}$~\cite{Steriade1993}. يمكننا تجميع هذه الأرجحة البطيئة من القطع المتاحة لدينا. لنأخذ مجموعة من عصبونات \nt{GLUT} بتغذية راجعة قوية على نفسها، أي الذاكرة من الفصل~\ref{sec:glutcomputer} التي لها حالتان مستقرتان، ولنضف تعبًا بطيئًا، كما في مولّد الإيقاع في الفصل~\ref{sec:muscles}:
\begin{empheq}[box=\keybox]{equation}
\tau\,\frac{\dd r}{\dd t} = -r + F\bigl(w\,r + I - a\bigr),
\qquad
\tau_a\,\frac{\dd a}{\dd t} = -a + b\,r .
\label{eq:updown}
\end{empheq}
تشتغل المجموعة وتعمل وتتعب؛ ويأكل التعبُ $a$ الحالةَ العليا، فتسقط المجموعة في الصمت؛ وفي الصمت يزول التعب، وتختفي الحالة الدنيا، فتشتغل المجموعة من جديد. فنحصل على المذبذب الاسترخائي نفسه الذي رأيناه في ضغط النوم، لكنه أسرع بنحو ثمانين ألف مرة.

\begin{figure}[t]
\centering
\begin{tikzpicture}[>=Stealth,x=2.9cm,y=1.0cm]
\begin{scope}[yshift=1.9cm]
\draw[->,gray!70!black] (0,0) -- (4.1,0);
\draw[->,gray!70!black] (0,0) -- (0,1.25) node[left,font=\small,black] {$r$};
\draw[thick,blue!60!black] plot file {figures/updown_sleep.dat};
\node[font=\scriptsize,anchor=west] at (4.15,0.5) {نوم: $b=5$};
\end{scope}
\draw[->,gray!70!black] (0,0) -- (4.1,0) node[right,font=\small,black] {$t$، ثانية};
\draw[->,gray!70!black] (0,0) -- (0,1.25) node[left,font=\small,black] {$r$};
\draw[thick,red!70!black] plot file {figures/updown_wake.dat};
\node[font=\scriptsize,anchor=west] at (4.15,0.5) {نهار: $b=1$};
\foreach \t in {0,1,2,3,4}{\draw[gray!60!black] (\t,-0.05) -- (\t,0.05); \node[below,font=\scriptsize] at (\t,0) {\t};}
\end{tikzpicture}
\caption{المجموعة نفسها~\eqref{eq:updown} في وضعين: $w=5$، $I=2$، $\theta=2.5$، $\tau=10\ms$، $\tau_a=0.3$ ثانية. مع تعب قوي ($b=5$، كما في النوم البطيء) تتأرجح المجموعة بين العمل والصمت بدور $1.1$ ثانية (قيمة النموذج؛ وفي القشرة الدور أطول من ثانية، وتحت التخدير نحو $3$ ثوانٍ). وإذا أضعفنا التعب ($b=1$)، وهذا ما يفعله الأسيتيل كولين، عملت المجموعة نفسها بانتظام.}
\label{fig:updown}
\end{figure}

فما الذي ينقل الشبكة من الأرجحة إلى العمل النهاري المنتظم؟ الأسيتيل كولين يثبّط في العصبونات التيارات البطيئة التي تصنع التعب~\cite{McCormick1992}. في نموذجنا، مع التعب القوي $b=5$ تتأرجح المجموعة بدور $1.1$ ثانية، وهو أسرع من الأرجحة المُقاسة لكنه من رتبتها، وتعمل نحو $35\%$ من الوقت إذا أُخذ المتوسط على عدد صحيح من الأدوار؛ ومع التعب الضعيف $b=1$ تعمل المجموعة نفسها بانتظام (الشكل~\ref{fig:updown}). سجل واحد، هو مستوى \nt{ACET}، يحوّل المذبذب إلى مضخم. وفوق الأرجحة البطيئة تسري في النوم ومضات أسرع من إيقاع $7\text{--}15$ ذبذبة في الثانية؛ تولّدها حلقة \foreignlanguage{english}{\nt{GLUT}--\nt{GABA}} بين القشرة وعقدة الترحيل الوسطى~\cite{SteriadeMcCormickSejnowski1993}، وهي قريبة الإيقاع الذي رأيناه في الفصل~\ref{sec:gaba}.

\section{الإعادات: تشغيل اليوم مسرَّعًا}

رأينا في الفصل~\ref{sec:acet} أن العصبونات التي عملت معًا نهارًا تطلق في النوم البطيء معًا من جديد وبالترتيب نفسه. ولنضف تفصيلًا هندسيًا واحدًا: الإعادات تجري \emph{مسرَّعة}. التسلسل الذي استغرق ثواني في النهار يُعاد في النوم في أعشار من الثانية~\cite{Nadasdy1999}: في الحُصين أسرع بنحو عشرين مرة~\cite{LeeWilson2002}، وفي القشرة بست أو سبع مرات~\cite{Euston2007}. ولا يُعاد في أي وقت اتفق: الومضات القصيرة للإعادات في منطقة واحدة تتزامن مع أرجحة العمل والصمت ومع الومضات السريعة للإيقاع في القشرة. وإذا عُزّز هذا التزامن اصطناعيًا تحسنت الذاكرة~\cite{Maingret2016}، وهذه علاقة سببية لا مجرد تزامن.

لماذا تشغّل الآلة يومها مسرَّعًا؟ يعرف المهندس صعوبة تسمى \emph{النسيان الكارثي}: الشبكة التي تتعلم الجديد متتابعًا، شيئًا بعد شيء، تفسد القديم. والمنقذ هو \emph{الخلط}: أن يُتعلم الجديد مخلوطًا بالقديم، على دفعات صغيرة، مرات كثيرة. وفرضية نظامي التعلم~\cite{McClelland1995} تقول إن الذاكرة السريعة تسجل اليوم كاملًا، ثم تعيد تشغيله في النوم للقشرة البطيئة شيئًا فشيئًا، مخلوطًا ومرات كثيرة. هذا هو الزوج نفسه ”مخزن مؤقت سريع، مستودع بطيء“ في الفصل~\ref{sec:acet}، والزوج نفسه من المتعلم السريع والمتعلم البطيء في الفصل~\ref{sec:motorlearning}. الإعادات وتزامنها مُقاسة؛ أما أن غرضها هو التعلم المخلوط للذاكرة البطيئة ففرضية جيدة الأساس.

\section{إعادة تطبيع الأوزان}

في النهار تقوّي قواعد هب من الفصل~\ref{sec:dofa} الوصلات في الغالب. ولو اقتصر الأمر على التقوية لنمت الأوزان، ونمت معها كلفة الطاقة وهبطت الحساسية: العصبون الذي كل مداخله قوية يكفّ عن التمييز بينها. وبحسب فرضية \emph{الاتزان الداخلي المشبكي}~\cite{TononiCirelli2014}، يلزم النوم، ضمن ما يلزم له، لإعادة الأوزان إلى مستوى العمل: تضعف الوصلات كلها بالنسبة نفسها، فتبقى \emph{نسبها}، أي ما تعلّمته الآلة. وهذا عند المهندس تطبيع، كما في القاعدة~\eqref{eq:oja}، غير أنه لا يجري أثناء العمل، بل في وقت منفصل.

والقياسات المباشرة لا تناقض ذلك: عند الفئران تكون مساحة التماس في مشابك القشرة بعد النوم أصغر بنحو $18\%$ مما بعد اليقظة، والتناقص متناسب مع الحجم، أما أكبر التماسات وأثبتها فلا تكاد تتغير~\cite{deVivo2017}. تحجيم الأوزان مُقاس؛ أما أنه المهمة الرئيسية للنوم ففرضية.

\section{النوم السريع: آلة مفصولة عن مداخلها ومخارجها}

في النوم السريع تعمل القشرة تقريبًا كما في النهار، لكن المدخل من أعضاء الحس مُضعَف، والمخرج إلى العضلات مطفأ بالتثبيط. هذا عند المهندس وضع مألوف: \emph{الاختبار على منصة}. الدارة تعمل بكامل قدرتها، لكن مداخلها موصولة بمولّد داخلي، ومخارجها مفصولة عن المشغّلات كي لا ينكسر شيء. والأحلام، بحسب إحدى الفرضيات، هي بالضبط هذا التشغيل الداخلي لنموذج العالم الذي جُمع خلال النهار؛ أما ما تفعله الشبكة حينئذ بالضبط وهل تتعلم، فغير معروف. المُقاس هنا هو فقط ما في الجدول~\ref{tab:sleepmodes}: أي قناة مشغّلة وأيها صامتة، وأن المخرج محجوب.

\section{الصيانة والنوم الموضعي}

ساد طويلًا أن الفراغات بين الخلايا تتسع في النوم فيُغسل الدماغ أسرع من فضلات الأيض~\cite{Xie2013}. لكن تجارب حديثة قاست التصريف بطريقة أخرى فوجدت العكس: في النوم يكون أبطأ~\cite{Miao2024}. السؤال مفتوح الآن، وسنتركه مفتوحًا.

وثمة حقيقة أوثق: النوم خاصية للشبكات الموضعية، لا للآلة كلها دفعة واحدة. إذا حُرم حيوان من النوم طويلًا، فإن مناطق منفردة من قشرته، والحيوان مستيقظ يتحرك، تسقط لحظات قصيرة في الصمت كما في النوم البطيء، ويخطئ الحيوان في تلك اللحظات أكثر~\cite{Vyazovskiy2011}. كل شبكة تتعب بنفسها وتتبدل بنفسها؛ ويتكون الوضع العام حين تنقل قنوات البث الشبكات كلها معًا.

\begin{keyblock}
\begin{proposition}[النوم مجموعةَ أوضاع]
\label{prop:sleep}
في النوم لا تنطفئ العصبونات: قنوات البث تنقل الآلة إلى أوضاع أخرى. ومتى ننام يقرره مُرحِّل بتخلّف~\eqref{eq:sleeppressure} يضبطه الإيقاع اليومي. في النوم البطيء تصير الشبكة مذبذبًا استرخائيًا~\eqref{eq:updown} وتعيد تشغيل اليوم مسرَّعًا؛ وفي النوم السريع تعمل مفصولةً عن المداخل والمخارج. الأوضاع والإعادات وتحجيم الأوزان مُقاسة؛ أما غرضها، أي التعلم المخلوط وإعادة التطبيع، ففرضيات جيدة الأساس.
\end{proposition}
\end{keyblock}

\section{الخلاصة}

\begin{table}[t]
\centering
\footnotesize
\setlength{\tabcolsep}{4pt}
\renewcommand{\arraystretch}{1.15}
\begin{tabular}{>{\raggedright}p{3.2cm}>{\raggedright}p{4.4cm}>{\raggedright\arraybackslash}p{5.2cm}}
\toprule
ماذا يحدث & كيف & ما المعروف \\
\midrule
تبدّل الأوضاع & السجلات \nt{ACET} و\nt{NORA} و\nt{SERO} (الجدول~\ref{tab:sleepmodes}) & مُقاس \\
متى ننام & المكثف $S$ ومُرحِّل بتخلّف~\eqref{eq:sleeppressure} & النموذج يصف التجارب جيدًا؛ الأدينوزين بوصفه $S$ فرضية \\
الأرجحة البطيئة & مجموعة بتغذية راجعة وتعب~\eqref{eq:updown} & الأرجحة مُقاسة؛ والنموذج أبسط مخطط \\
تشغيل اليوم مسرَّعًا & إعادات أسرع بـ$6\text{--}20$ مرة، متزامنة مع الأرجحة & مُقاس؛ وتعزيز التزامن يحسّن الذاكرة \\
لماذا الإعادات & التعلم المخلوط للذاكرة البطيئة & فرضية جيدة الأساس \\
إعادة تطبيع الأوزان & تحجيم الوصلات في النوم & تناقص التماسات مُقاس؛ وكونه الدور الرئيسي للنوم فرضية \\
النوم السريع & العمل مع مداخل ومخرج مفصولة & الوضع مُقاس؛ والغرض غير معروف \\
التصريف & اتساع الفراغات بين الخلايا & نتائج متناقضة \\
النوم الموضعي & مناطق منفردة تسقط في الصمت & مُقاس \\
\bottomrule
\end{tabular}
\caption{ماذا تفعل العصبونات في النوم.}
\label{tab:sleep}
\end{table}

يلخص الجدول~\ref{tab:sleep} الفصل. تبيّن أن النوم ليس توقفًا في عمل الآلة، بل صيانتها أثناء العمل: تبديل الأوضاع بقنوات البث نفسها، ومولّد من القطع نفسها التي يُبنى منها إيقاع الخطوات، وتشغيل اليوم مسرَّعًا للذاكرة البطيئة، وإعادة تطبيع الأوزان. في النهار تكتب الآلة، وفي الليل تعيد كتابة ما كُتب وترتّبه.

\section{تمارين}

\begin{exercise}[\lvlA]
\label{ex:20:1}
\emph{مُرحِّل بلا إيقاع يومي.} لتكن العتبتان ثابتتين: $H=0.67$، $L=0.17$ (متوسطا عتبتي نموذج الشكل~\ref{fig:twoprocess}). استخرج من~\eqref{eq:sleeppressure} مدتي اليقظة والنوم ودور المذبذب عند $\tau_r=18.2$ ساعة و$\tau_d=4.2$ ساعة. لماذا يصل النموذج ذو العتبات المتأرجحة مع ذلك إلى $16.1+8.0=24$ ساعة؟ وإلى أي جهاز في الفصل~\ref{sec:chemoutput} يعود هذا الأثر؟
\end{exercise}

\begin{exercise}[\lvlA]
\label{ex:20:2}
\emph{دَين النوم.} النوم الذي يبدأ عند ضغط $S_0$ يدوم $t_s=\tau_d\ln(S_0/L)$، $\tau_d=4.2$ ساعة، و$L$ العتبة الدنيا عند الاستيقاظ. (أ)~بعد $40$ ساعة بلا نوم $S_0=0.90$، ودام النوم $8.8$ ساعة. كم كانت $L$؟ وكم كان النوم سيدوم مع $L\approx0.092$ المعتادة؟ (ب)~خلال الليل تنقص مساحة التماسات المشبكية $18\%$. بكم يجب أن تنمو الأوزان خلال النهار؟
\end{exercise}

\begin{exercise}[\lvlB]
\label{ex:20:3}
\emph{تصميم العتبات.} اختر عتبتين ثابتتين $H$ و$L$ يعطي عندهما المُرحِّل~\eqref{eq:sleeppressure} مع $\tau_r=18.2$ ساعة و$\tau_d=4.2$ ساعة $16$ ساعة يقظة و$8$ ساعات نوم بالضبط. بمَ تكون هذه الآلة أسوأ من الآلة ذات العتبات المتأرجحة إذا أُوقظت ذات ليلة بعد $4$ ساعات؟
\end{exercise}

\begin{exercise}[\lvlB]
\label{ex:20:4}
\emph{دور الأرجحة البطيئة.} في النموذج~\eqref{eq:updown} $F(x)=1/\bigl(1+e^{-(x-\theta)/k}\bigr)$، $w=5$، $I=2$، $\theta=2.5$، $k=0.25$، $b=5$، $\tau\ll\tau_a=0.3$ ثانية. (أ)~أوجد نقطتي الانعطاف $a_{\text{hi}}$ و$a_{\text{lo}}$ للمنحنى $r=F(wr+I-a)$، حيث يختفي العمل والصمت. (ب)~باعتبار $r\approx1$ في العمل و$r\approx0$ في الصمت، أوجد الدور ونسبة العمل.
\end{exercise}

\begin{exercise}[\lvlB]
\label{ex:20:5}
\emph{متى تنطفئ الأرجحة.} في نموذج التمرين~\ref{ex:20:4} تقع النقطة الثابتة عند تقاطع المنحنى $a=wr+I-F^{-1}(r)$ مع المستقيم $a=br$. تستمر التذبذبات ما دام التقاطع على الفرع الأوسط، بين نقطتي الانعطاف. عند أي $b$ يكون ذلك؟ وكيف يحوّل \nt{ACET} المذبذب إلى مضخم بتغيير $b$؟
\end{exercise}

\begin{exercise}[\lvlB]
\label{ex:20:6}
\emph{نمذجة: دَين أم طور.} في \texttt{sleep\_models.py} خذ الاستيقاظ الأول بعد $48$ ساعة وأبقِ الآلة مستيقظة $D=24$ و$32$ و$40$ و$48$ و$64$ ساعة (\texttt{stay\_awake}، \texttt{days=8}). كم يدوم نوم التعافي عند كل $D$؟ وما الذي يحدد مدته؟
\end{exercise}

\begin{exercise}[\lvlB]
\label{ex:20:7}
\emph{نمذجة: النسيان والخلط.} عصبون خطي $y=\mathbf w\cdot\mathbf x$، $n=40$، يتعلم بالقاعدة $\mathbf w\leftarrow\mathbf w-0.5\,(y-y^*)\,\mathbf x$. المهمتان A وB، لكل منهما $10$ أنماط، $x_j\sim\mathcal N(0,1/n)$، $y^*=\pm1$. ما جذر متوسط مربع الخطأ على A بعد $200$ حقبة من A ثم $200$ حقبة من B؟ وبعد $200$ حقبة من A وB مخلوطتين؟
\end{exercise}

\begin{exercise}[\lvlC]
\label{ex:20:8}
\emph{بحث: الخلط أم إعادة التطبيع؟} عصبون التمرين~\ref{ex:20:7} مع عتبة؛ وإجراءان ليليان: (i)~تُضرب كل الأوزان في $0.82$؛ (ii)~تُعاد الأنماط القديمة مخلوطة بالجديدة. ماذا يفعل كل منهما بنسب الأوزان وباستجابات العصبون؟ وكيف يمكن التمييز بين الفرضيتين بأحجام المشابك بعد النوم؟
\end{exercise}
